\section{Related Work} \label{sec:relatedwork}
\subsection{Domain-generalizable person re-identification} \label{subsec:DG-ReID}
DG-ReID methods often seek domain-invariant embeddings that transfer to unseen target cameras without fine-tuning. Instance normalization and hybrids such as IBN-Net \cite{pan2018two} and style-normalization variants SNR \cite{jin2020style} attempt to filter out domain-specific style variations by manipulating feature statistics. Similarly, domain-alignment and meta-learning formulations \cite{song2019generalizable,choi2021meta} explicitly optimize for feature representations that minimize divergence across seen source domains. Balancing Alignment and Uniformity (BAU) \cite{cho2024generalizable} instantiates the alignment-uniformity framework of \citet{wang2020understanding} for multi-source ReID, explicitly balancing alignment and uniformity on originals and augmentations while adding domain-specific uniformity to exploit augmentation without collapsing out-of-distribution structure. Collectively, these approaches emphasize suppressing domain-specific factors while sharing a common retrieval convention: symmetric distances on Euclidean (or cosine-normalized) features, which is natural when a single bidirectional gallery score is the operational target.

% In constrast, our Finsler-geometric framework operates as a \textit{property-preserving} model, structurally encoding domain shifts as asymmetrical traversal costs rather than erasing the underlying domain physics.

\subsection{Retrieval geometry} \label{subsec:retrieval_geometry}
Batch-hard triplet and proxy losses for ReID are almost universally built from symmetric distances \cite{hermans2017defense}. Cross-modality settings (e.g., RGB–infrared) motivate directional or modality-specific heads and losses \cite{wu2017rgb}, where the asymmetry arises from sensor heterogeneity. A separate line replaces Euclidean feature space with non-Euclidean intrinsic geometry: hyperbolic embeddings \cite{khrulkov2020hyperbolic} and hyperspherical manifold methods \cite{fan2019spherereid} modify the intrinsic geometry of the representation space itself for metric learning and retrieval. A fourth design axis keeps a standard Euclidean feature map and injects direction dependence in the pairwise score rather than in the geodesics of the embedding space; Finsler-MDS \cite{dages2025finsler} is a recent example of this template (see Sec.~\ref{subsec:asymmetric_finsler_distance}).

\subsection{Feature disentanglement} \label{subsec:feature_disentanglement}
Disentangling identity from style or domain factors appears in both generative and discriminative pipelines. Generative approaches such as DG-Net \cite{zheng2019joint} and MUNIT \cite{huang2018multimodal} use adversarial image synthesis to swap identity and style codes. While effective at isolating appearance factors, GAN training is resource-intensive and the learned subspaces do not typically yield strict mutual statistical independence.

Discriminative approaches apply dedicated identity and nuisance losses within a shared backbone \cite{jin2020style}, avoiding image synthesis. In both families, orthogonality constraints are inductive biases and do not guarantee strict factorization in practice. In joint training with strong domain-invariance terms, gradients from nuisance-focused objectives can oppose invariance pressures on shared representations. This represents a tension in multi-task metric learning that is distinct from the theoretical question of whether disentanglement is desirable.

% Explicitly factoring latent spaces into independent \textit{identity} and \textit{domain} subspaces has driven significant advancements in generative and discriminative CV tasks. Pioneering architectures like DG-Net \cite{zheng2019joint} and MUNIT \cite{huang2018multimodal} utilize adversarial training to swap identity and style codes, effectively isolating the identity manifold. However, these disentanglement strategies often rely on computationally intensive generative adversarial networks and struggle to guarantee strict orthogonality in the latent space. Our architectural contribution grounds feature disentanglement within a rigorous geometric trajectory. Rather than relying on adversarial generation, we explicitly map identity features to the invariant structural components of the Finsler manifold, while confining domain variations (e.g., extrinsic camera parameters) to the asymmetric, tangent-space routing functions. This ensures mathematically guaranteed separation of the identity and domain subspaces, maintaining discriminative power without the instability of adversarial training. 

\begin{figure*}[t]
  \centering
  \includegraphics[width=\linewidth]{fig/fig_finslerReID_arch.pdf}
  \caption{\textbf{Finsler BAU framework.} Two image streams $\mathbf{x}$ and $\hat{\mathbf{x}}$ depict the same identity observed under different camera conditions; both pass through the shared backbone and drift head. \emph{Right:} loss terms acting on each component. $\mathcal{L}_{\mathrm{dcc}}$ (eq.~\ref{eq:loss_dcc}) penalises L2 dispersion between drift vectors $\omega_i, \omega_j$ for same-identity, different-camera pairs. $\mathcal{L}_\omega$ (eq.~\ref{eq:omega_regularization}) enforces the Randers positivity constraint on \emph{all} drift vectors. The bidirectional Finsler triplet $\mathcal{L}_{\mathrm{tri}}^{\mathrm{bi}}$ (eq.~\ref{eq:bidirectional_triplet}) mines hard pairs in both query directions on the full structured embedding $[\mathbf{z}^{\mathrm{id}};\hat{\mathbf{z}}^{\omega}]$.} 
  \label{fig:architecture}
\end{figure*}

\subsection{Randers-Finsler geometry} \label{subsec:randers-finsler_geometry}
\paragraph{Classical Randers metrics.} A Randers metric combines a Riemannian quadratic piece and a linear drift in the tangent direction: $F_x(u)=\sqrt{u^\top M(x)\,u}+\omega(x)^\top u$, with positivity ensured by $\|\omega(x)\|_{M(x)^{-1}}<1$ \cite{randers1941geometry,bao2000introduction}. The Riemannian (symmetric) case is $\omega\equiv 0$; the drift term $\omega$ breaks symmetry by making traversal cost depend on direction.

\paragraph{Finsler multi-dimensional scaling (Finsler-MDS).} Finsler MDS \cite{dages2025finsler} embeds asymmetric dissimilarities into a flat canonical Randers space on $\mathbb{R}^m$: a uniform drift $\omega\in\mathbb{R}^m$ with $\alpha\equiv\|\omega\|_2<1$ defines
\begin{equation}
    F_x^{C}(u)=\|u\|_2+\omega^\top u ,
    \label{eq:finsler_norm}
\end{equation}
with straight Euclidean segments as geodesics but direction-dependent traversal length. The canonical Randers distance is then
\begin{equation}
    d_{F^{C}}(x,y)=\|y-x\|_2+\omega^\top (y-x) ,
    \label{eq:finsler_dist_canonical}
\end{equation}
so $d_{F^{C}}(y,x)\neq d_{F^{C}}(x,y)$ when $\omega\neq 0$, and $\omega=0$ recovers Euclidean distance \cite{dages2025finsler}. The additive form, Euclidean chord plus a linear functional of the displacement, yields a closed-form asymmetric cost amenable to MDS stress minimization and graph embedding. Extensions to learned or position-dependent drift fields must satisfy integrand-level regularity conditions to define genuine Finsler metrics \cite{bao2000introduction}. \\

\subsection{Asymmetric domain shift.} \label{subsec:assymmetric_domain_shift}
Cross-modal and viewpoint-heterogeneous ReID violate the idealization that ``domain distance'' is symmetric under query reversal. Visible-infrared retrieval is inherently directional because sensor physics induce irreversible information loss~\cite{wu2017rgb}: SYSU-MM01 reports RGB$\to$IR and IR$\to$RGB separately with nontrivial Rank-1 gaps between directions. Aerial-to-ground benchmarks on AG-ReID.v2 likewise use directional probe-gallery splits, and performance depends on which view is queried~\cite{nguyen2024agreidv2}. For generic degradations, the data processing inequality makes the asymmetry formal: a lossy channel cannot be inverted, so clean$\to$corrupted information loss typically dominates the reverse~\cite{cover2006elements}. These observations motivate treating probe$\to$gallery scoring as directed, which we instantiate in Sec.~\ref{sec:methodology}. \citet{chen2022benchmarks} make this direction asymmetry directly measurable: their Table~2 reports Rank-1 of $30.5\%$ when the query is corrupted versus $77.2\%$ when the gallery is corrupted under a ResNet-50 baseline on Market-1501, a $46$-point gap driven by which side of the retrieval pair is degraded. Our balanced bidirectional evaluation (Sec.~\ref{subsec:controlled_corruption}) controls for this protocol effect by equalizing query count, gallery size, and correct-match count across both directions, and confirms the same direction: clean-probe retrieval against a corrupted gallery is reliably easier than the reverse at all tested severity levels.

% Cross-modal and viewpoint-heterogeneous ReID violate the idealization that ``domain distance'' is symmetric under query reversal. Visible-infrared retrieval is inherently directional because sensor physics induce irreversible information loss~\cite{wu2017rgb}; SYSU-MM01 reports RGB$\to$IR and IR$\to$RGB separately with nontrivial Rank-1 gaps between directions. Aerial-to-ground benchmarks on AG-ReID.v2 likewise use directional probe-gallery splits, and performance depends on which view is queried~\cite{nguyen2024agreidv2}. For generic degradations, the data processing inequality makes the asymmetry formal: a lossy channel cannot be inverted, so clean$\to$corrupted information loss typically dominates the reverse~\cite{cover2006elements}. These observations motivate treating probe$\to$gallery scoring as directed, which we instantiate in Sec.~\ref{sec:methodology}.

The preceding threads share a gap: none simultaneously (a) operates within a standard DG-ReID training recipe, (b) keeps the backbone representation Euclidean, and (c) introduces direction-dependent retrieval costs through an outer asymmetric score rather than through intrinsic embedding geometry or modality-specific fusion. We build on BAU \cite{cho2024generalizable} as the invariance backbone and adopt the displacement-linear template of $d_{F^{C}}$ for the scoring layer (see Sec.~\ref{subsec:asymmetric_finsler_distance}).